% Clase 22 --- Las tres impedancias elementales y su desfasaje (§3.3).
% La tabla del texto da los tres valores; lo que la tabla no muestra es qué
% significan: cada j equivale a un cuarto de período de corrimiento entre la
% tensión y la corriente. Por eso cada panel lleva su par de curvas.
\begin{center}
\begin{tikzpicture}[scale=1]

  \providecommand{\parcurvas}[1]{%
    \begin{axis}[
      fisfig, width=5.0cm, height=3.2cm,
      axis lines=middle, grid=none,
      xlabel={}, ylabel={},
      xmin=0, xmax=6.5, ymin=-1.5, ymax=1.5,
      xtick=\empty, ytick=\empty,
      every axis plot/.append style={line width=0.9pt},
    ]
      \addplot[figblue, smooth, samples=140, domain=0:6.4] {cos(deg(x))};
      \addplot[figred, smooth, samples=140, domain=0:6.4]
        {cos(deg(x - #1))};
    \end{axis}}

  % =============== (a) resistencia ===============
  \begin{scope}
    \draw (0.2,2.6) to[R, l=$R$] (2.6,2.6);
    \parcurvas{0}
    \node[figlbl, anchor=north, align=center] at (1.4,-1.35)
      {(a) $\hat Z_R = R$\\[-0.1em] en fase};
  \end{scope}

  % =============== (b) bobina ===============
  \begin{scope}[xshift=5.0cm]
    \draw (0.2,2.6) to[L, l=$L$] (2.6,2.6);
    \parcurvas{-90}
    \node[figlbl, anchor=north, align=center] at (1.4,-1.35)
      {(b) $\hat Z_L = j\omega L$\\[-0.1em] $I$ atrasa $90^\circ$};
  \end{scope}

  % =============== (c) condensador ===============
  \begin{scope}[xshift=10.0cm]
    \draw (0.2,2.6) to[C, l=$C$] (2.6,2.6);
    \parcurvas{90}
    \node[figlbl, anchor=north, align=center] at (1.4,-1.35)
      {(c) $\hat Z_C = \dfrac{1}{j\omega C}$\\[-0.1em] $I$ adelanta $90^\circ$};
  \end{scope}

  \node[figlbl, anchor=north] at (6.4,-3.05)
    {\textcolor{figblue}{---} tensión \qquad
     \textcolor{figred}{---} corriente};

\end{tikzpicture}\\[0.5em]
{\small\itshape Las tres impedancias elementales dicen dos cosas a la vez, y
conviene leerlas separadas. El módulo dice cuánto ``cuesta'' hacer pasar
corriente: constante en la resistencia, creciente con la frecuencia en la bobina
---que se opone a los cambios rápidos--- y decreciente en el condensador, que a
frecuencia alta apenas alcanza a cargarse y casi no estorba. El factor $j$ dice
lo otro: un cuarto de período de corrimiento entre tensión y corriente. En la
resistencia no hay $j$ y las dos curvas van juntas; en la bobina la corriente
queda atrasada y en el condensador adelantada, y esa oposición entre ambas es
justamente la que permite que se cancelen en resonancia. Una vez que cada
elemento tiene su número complejo, los circuitos se combinan con las reglas de
siempre ---en serie se suman las impedancias, en paralelo se suman los
inversos---, sólo que ahora la aritmética es compleja. De $\hat Z_{\text{eq}}$
se sacan el módulo (amplitud) y el argumento (desfasaje).}
\end{center}
